\documentclass[journal]{IEEEtran}

\usepackage{graphicx}
\usepackage{booktabs}
\usepackage{amsmath}
\usepackage{url}
\usepackage{cite}
\usepackage{array}
\usepackage{multirow}
\usepackage{balance}

\begin{document}

\title{A Privacy-Preserving Provenance Plane for Reconstructable Fast Healthcare Interoperability Resources to Large Language Model Clinical Artificial Intelligence Transactions}

\author{Satya~Venkata~Ranga~Janaki~Sriram~Mentey%
\thanks{S.~V.~R.~J.~S.~Mentey is the corresponding author.
Address: 8601 Anderson Mill Rd, Apt~722, Austin, TX~78729, USA
(e-mail: srirammentey@ieee.org; ORCID: 0009-0007-2681-006X).}}

\markboth{IEEE Journal of Biomedical and Health Informatics}%
{Mentey: Provenance Plane for FHIR-to-LLM Clinical AI Transactions}

\maketitle

\begin{abstract}
Clinical artificial intelligence systems that consume electronic health record
data and emit reviewer-facing summaries rarely preserve the inputs,
transformations, model configuration, evidence, safeguards, and human
disposition needed to reconstruct a transaction, while storing raw clinical
payloads or hidden model reasoning in an immutable log is unsafe. This paper
presents Curie Audit Plane, a provenance method for one bounded synthetic
Fast Healthcare Interoperability Resources Release~4 to large language model
transaction. The plane separates protected clinical content from an
append-only audit chain, records typed events with content digests, and
verifies hash chaining, Merkle inclusion, and Ed25519 signatures. Structured
model output is constrained to rationale, evidence references, uncertainty,
assumptions, and missing data. On synthetic fixtures, independently verified
audit reconstruction completeness is $20/20$ required fields, all $19$ labeled
tamper cases are detected, and access-control probes pass $11/11$ allowed or
denied outcomes. Capture overhead is reported against an identical unrecorded
clinical workflow, not a thin log shortcut. Significance for biomedical
research: the method shows how a FHIR-to-LLM handoff can be reconstructed and
independently verified without placing hidden chain-of-thought or raw
protected health information in an immutable store.
\end{abstract}

\begin{IEEEkeywords}
Biomedical informatics, Fast Healthcare Interoperability Resources,
clinical artificial intelligence, provenance, auditability, integrity,
health information systems.
\end{IEEEkeywords}

\section{Introduction}
\label{sec:intro}

Large language models are being attached to electronic health record workflows
for summarization and review support. Healthcare organizations still need a
narrower answer: given one completed transaction, can an investigator
reconstruct what the model was shown, how it was configured, which evidence
and safeguards were applied, and what a human reviewer decided, without
placing raw protected health information or hidden reasoning into an immutable
store?

This paper treats that question as a biomedical and health-informatics
provenance problem, not as generic software logging. The prototype is
deliberately narrow: one synthetic Fast Healthcare Interoperability Resources
(FHIR) Release~4 encounter, one structured clinical-summary task, and a human
ACCEPT/MODIFY/REJECT disposition.

The frozen research questions are: (i)~reconstructability of required
provenance fields; (ii)~detection of mutation, omission, reordering, broken
references, and invalid proofs; (iii)~latency and storage overhead of complete
capture versus an identical unrecorded clinical workflow; (iv)~console exposure
of evidence, uncertainty, guardrail results, and disposition as a system
demonstration (\texttt{SCRIPTED\_PROXY}), not a blinded reviewer study; and
(v)~replay fidelity for a deterministic stub versus a hosted nondeterministic
model.

\section{Related Work}
\label{sec:related}

Healthcare provenance and audit have a long FHIR-native literature, including
W3C PROV mappings and FHIR \texttt{Provenance}/\texttt{AuditEvent} resources
used to record agents, entities, and access~\cite{hl7provenance,hl7auditevent,provdm}.
Those resources are necessary interoperability projections; they are not, by
themselves, an internal event chain with independent cryptographic
verification.

Clinical decision-support logging and electronic health record audit trails
typically record who accessed what, not the full artificial-intelligence
transaction: input manifests, transformation records, prompt and model
identity, retrieved evidence digests, guardrail outcomes, and the human
handoff. Model-provider traces and application JSON logs can retain payloads
that should not be immutable, or they omit the links required for
reconstruction.

This work is closest to tamper-evident logging and hash-chained
audit~\cite{schneier1998logs,crosby2009tamper}, combined with a clinical-data
boundary. It does not replace deterministic alert governance, FHIR access
control, or the large language model provider. Related comparisons in the
evaluation are application JSON logging, hash-only records without signatures,
and FHIR-only projection.

\section{Method}
\label{sec:method}

\subsection{System Boundary}
\label{sec:boundary}

Curie Audit Plane observes a synthetic FHIR-to-large-language-model workflow.
Protected payloads (assembled context, model request and response,
reviewer-modified text, and free-form comments) are stored in an authorized
content store. The immutable audit store holds metadata, opaque references,
schemas, and digests. Reviewer comments on the chain are limited to a
controlled category; detailed text is protected. Provider exception text is
not stored or fingerprinted. Patient and resource identifiers in immutable
metadata are opaque tokens; a protected identity map is required for
authorized reconstruction.

\subsection{Event Model and Integrity}
\label{sec:events}

A successful transaction emits typed events from input manifestation through
human action and an integrity-proof commitment. Events are canonically hashed,
chained, batched with a Merkle root, and signed. Verification reports
\texttt{VERIFIED}, \texttt{INCOMPLETE}, \texttt{TAMPERED}, or \texttt{FAILED}.
FHIR Release~4 \texttt{Provenance} and \texttt{AuditEvent} resources are
documented projections of the internal model, not a claimed Implementation
Guide profile.

\subsection{Structured Output and Guardrails}
\label{sec:guardrails}

Model output must validate as structured rationale. Guardrails record schema
validity, evidence-reference coverage, uncertainty presence, a synthetic
identifier scan, and policy action. Blocking results require an explicit
override policy before ACCEPT. Hidden chain-of-thought fields are rejected.

\subsection{Evaluation Protocol}
\label{sec:eval}

The reproducible command is \texttt{curie-audit-plane evaluate}. Reports use
schema \texttt{curie-evaluation.v1.1} and include experiment metadata (git
commit, dirty-tree flag, fixture alias and digest, provider, model, prompt
version, decoding parameters, endpoint class, seed, and command template)
without absolute local paths as publication identifiers.

Headline audit reconstruction completeness (\texttt{independently\_verified\_arc})
is required-field presence after reloading persisted records, counted only when
independent verification succeeds. Field-presence completeness on the
in-memory object is reported separately. Tamper detection uses $19$ labeled
mutation cases. Replay is classified as \texttt{EXACT\_MATCH},
\texttt{EQUIVALENT}, or \texttt{DIVERGENT}. Overhead compares the complete
plane with an identical unrecorded clinical workflow (load, transform,
context, retrieve, complete, guardrails, JSONL) on isolated stores with warmup
and three timed repeats:
\begin{align}
\mathrm{latency\_ratio} &= (T_{\mathrm{plane}}-T_{\mathrm{no\_audit}})/T_{\mathrm{no\_audit}}, \\
\mathrm{storage\_ratio} &= (B_{\mathrm{plane}}-B_{\mathrm{no\_audit}})/B_{\mathrm{no\_audit}}.
\end{align}
Rate metrics use Wilson intervals; latency ratios use paired differences.
Ablations omit input manifests, transformations, model and prompt metadata,
evidence references, cryptographic proofs, or human provenance.
Access-control probes record allowed and denied HTTP outcomes for reviewer,
investigator, admin, output, content, export, missing transaction, and global
list scope. Application JSONL and hash-only JSONL are separately instrumented
from the unrecorded workflow; the FHIR projection is built from the source
bundle. These recorders are not independently shipped products. The complete
plane is scored from the signed audit chain.

The default $50$-encounter cohort rewrites one synthetic fixture. It is a
single-workflow protocol demonstration, not a population sample.
Optional Synthea slices may run from approved roots when dumps exist, but they
are excluded from measured evidence because generator version, seed, module set,
and CLI are not pinned~\cite{walonoski2018synthea}.

\subsection{Threat Model}
\label{sec:threat}

Adversaries may mutate, delete, or reorder events; substitute proofs or keys;
present unsupported findings; or attempt to place protected health information
in reviewer comments, provider errors, research exports, or endpoint
configuration. The prototype assumes synthetic data, local keys, and an
approved loopback model endpoint. It does not claim protection against
compromised signing keys or a malicious protected-content store operator.

\section{Results}
\label{sec:results}

All numbers below are from the frozen stub evaluation campaign in
\texttt{evaluation-results/} (schema \texttt{curie-evaluation.v1.1}).
Software-quality checks at the time of this draft were $143$ Python tests,
$94.59\%$ statement coverage, and a clean Ruff pass; those are not clinical
performance.

\subsection{Reconstructability and Tamper Detection}
Independently verified audit reconstruction completeness is $20/20$ required
fields on the clean fixture transaction. Field-presence completeness on the
in-memory object is also $20/20$. Required-event completeness is $11/11$ for
a successful transaction. All $19$ labeled tamper cases are detected, and the
false-tamper rate over three clean runs is $0$. On the cloned $50$-encounter
cohort, independently verified completeness remains $1.00$ with Wilson $95\%$
interval $[0.93, 1.00]$ ($n=50$) after reloading persisted records.

\begin{table}[!t]
\centering
\caption{Baseline comparison on the synthetic fixture (complete plane versus
separately instrumented unrecorded-workflow and source-bundle recorders).}
\label{tab:baselines}
\begin{tabular}{lccc}
\toprule
Recorder & ARC & Tamper detection & Queryability \\
\midrule
Application JSONL & $0.10$ & $0.00$ & $0.00$ \\
Hash-only JSONL & $0.15$ & $1.00$ & $0.00$ \\
FHIR R4 projection & $0.20$ & $0.00$ & $0.25$ \\
Complete plane & $1.00$ & $1.00$ & $1.00$ \\
\bottomrule
\end{tabular}
\end{table}

The complete plane reconstructs more required fields and detects the standard
model-event mutation that application logs and FHIR projections do not detect
(Table~\ref{tab:baselines}). Hash-only logging detects digest mismatch without
providing chain or signature guarantees.

\subsection{Ablations}
Omitting cryptographic proofs drops reconstructability from $1.00$ to $0.75$
($\Delta=0.25$). Omitting human provenance drops it to $0.85$
($\Delta=0.15$). Omitting model and prompt metadata drops it to $0.90$.
Omitting input manifests, transformations, or evidence references each drops
one required field ($0.95$). These are field-presence ablations of recorded
provenance classes, not clinical-efficacy ablations
(Table~\ref{tab:ablation}).

\begin{table}[!t]
\centering
\caption{Reconstructability after omitting a provenance class.}
\label{tab:ablation}
\begin{tabular}{lcc}
\toprule
Condition & ARC & $\Delta$ from full \\
\midrule
Full recorded provenance & $1.00$ & $0.00$ \\
Omit input manifests & $0.95$ & $0.05$ \\
Omit transformations & $0.95$ & $0.05$ \\
Omit model metadata & $0.90$ & $0.10$ \\
Omit evidence references & $0.95$ & $0.05$ \\
Omit cryptographic proofs & $0.75$ & $0.25$ \\
Omit human provenance & $0.85$ & $0.15$ \\
\bottomrule
\end{tabular}
\end{table}

\subsection{Access Control}
Eleven role probes passed with observed HTTP status equal to the predeclared
expectation ($11/11$; Wilson $95\%$ interval $[0.74, 1.00]$): reviewer read and
output allowed; reviewer export and investigator output or content denied;
investigator verify and export allowed; admin content allowed;
unauthenticated list denied ($401$); missing transaction reported as missing
($404$); global list allowed.

\subsection{Overhead and Replay}
Mean plane latency is $14.78$\,ms versus $9.15$\,ms for the identical
unrecorded workflow (paired latency ratio $0.615$, $95\%$ interval
$[0.594, 0.637]$, $n=3$ after one warmup). Storage ratio versus that workflow is
$6.31$ ($81\,885$ versus $11\,209$ bytes). These values do not support a claim
that audit capture is under $15\%$ of an already-complete clinical
application. Verification of the clean transaction is $1.07$\,ms.
Deterministic-stub same-prompt replay is exact; prompt-substitution replay is
\texttt{DIVERGENT} by construction. Hosted-model same-prompt replay is at best
\texttt{EQUIVALENT} and is often \texttt{DIVERGENT}.

The $16$-arm scenario matrix exercises ACCEPT/MODIFY/REJECT, forced and
natural WARN/BLOCK, provider failure, sparse encounter, Synthea slices,
MODIFY evidence, prompt-substitution replay, and access audit. It is workflow
coverage, not a population sample. The two Synthea slice arms are unpinned
demonstrations and are excluded from measured tables. Console reconstruction is
\texttt{SCRIPTED\_PROXY} only.

\section{Limitations}
\label{sec:limitations}

This is a single-workflow protocol demonstration on synthetic FHIR fixtures.
The default cohort rewrites one fixture; it is not $50$ independent patients
and does not support population-style inference. FHIR projections are mapping
rules, not Implementation Guide profile validation. Synthea generator version
and CLI parameters are \texttt{NOT\_PINNED}; optional slice arms are excluded
from measured evidence. A DOI for the public archive is not assigned in this draft.

\section{Ethics, Funding, and Disclosure}
\label{sec:ethics}

No human subjects and no identifiable clinical records were used. The reported
evaluation uses synthetic templates and optional external Synthea
bundles~\cite{walonoski2018synthea}. Therefore no institutional review board
protocol was required. The author reports no specific funding and no conflict
of interest. Portions of software and documentation were drafted with
assistance from coding agents; evaluation metrics are produced by the
versioned Python harness. The author remains responsible for claims, code, and
the decision to submit.

\section{Conclusion}
\label{sec:conclusion}

A privacy-preserving provenance plane can reconstruct and independently verify
a bounded synthetic FHIR-to-large-language-model transaction, detect a defined
mutation suite, and keep hidden reasoning and free-form reviewer text out of
the immutable chain. Remaining work is experimental scale on independently
generated synthetic patients, hosted-model replay characterization, and IEEE
Author Portal packaging (PDF Checker, signed consent if required).

\bibliographystyle{IEEEtran}
\bibliography{references}

\end{document}
